\section*{Danh mục các thuật ngữ}
\addcontentsline{toc}{section}{Danh mục các thuật ngữ}

% Điều chỉnh khoảng cách giữa các dòng trong bảng cho thoáng hơn
\renewcommand{\arraystretch}{1.3} 

\begin{longtable}{p{6.5cm} p{8.5cm}}
%─────────────────────────────────────────────────────────────────────────────
% TIÊU ĐỀ BẢNG - TRANG ĐẦU TIÊN
%─────────────────────────────────────────────────────────────────────────────
\rowcolor{headerblue!25}
\multicolumn{2}{l}{\bfseries\color{headerblue}
    \quad Thuật ngữ tiếng Việt / Tiếng Anh} \\
\endfirsthead

%─────────────────────────────────────────────────────────────────────────────
% TIÊU ĐỀ BẢNG - CÁC TRANG TIẾP THEO
%─────────────────────────────────────────────────────────────────────────────
\rowcolor{headerblue!25}
\multicolumn{2}{l}{\bfseries\color{headerblue}
    \quad Thuật ngữ tiếng Việt / Tiếng Anh (tiếp theo)} \\
\endhead

\hline
\multicolumn{2}{|r|}{\textit{Tiếp tục ở trang sau...}} \\
\hline
\endfoot

\endlastfoot

%─────────────────────────────────────────────────────────────────────────────
% NỘI DUNG DANH MỤC
%─────────────────────────────────────────────────────────────────────────────
\textbf{Biểu cảm nguyên mẫu} & Prototypic emotional expression \\
\textbf{Bộ lọc xác suất} & Probability filter \\
\textbf{Bộ lấy mẫu cân bằng} & Balanced sampler \\
\textbf{Bộ mã hóa văn bản} & Text encoder \\
\textbf{Bộ thích ứng thời gian} & Temporal adapter \\
\textbf{Bộ tiêu chuẩn đánh giá tĩnh} & Static benchmark \\
\textbf{Cấu trúc mạng phần dư} & Residual network \\
\textbf{Chỉ số đo lường} & Evaluation metric \\
\textbf{Chuyển đổi đặc trưng} & Feature projection \\
\textbf{Cơ chế chú ý chéo} & Cross-attention \\
\textbf{Cơ chế phòng vệ} & Defense mechanism \\
\textbf{Cơ chế rớt kênh ngẫu nhiên} & Random channel-dropping \\
\textbf{Đặc trưng biểu cảm đã chọn lọc} & Filtered expression feature \\
\textbf{Đặc trưng biểu cảm thô} & Raw expression feature \\
\textbf{Đặc trưng hình học} & Geometric feature \\
\textbf{Đặc trưng học sâu} & Learned feature \\
\textbf{Đặc trưng khuôn mặt} & Face feature \\
\textbf{Đặc trưng ngoại hình} & Appearance feature \\
\textbf{Đặc trưng thủ công} & Hand-crafted feature \\
\textbf{Độ chính xác trung bình} & Mean accuracy \\
\textbf{Độ chính xác tổng thể} & Overall accuracy \\
\textbf{Động lực học thời gian} & Temporal dynamics \\
\textbf{Đơn vị hành động cơ mặt} & Action unit \\
\textbf{Giao thức thực nghiệm} & Experimental protocol \\
\textbf{Gợi ý ngôn ngữ - hình ảnh} & Vision-language prompt \\
\textbf{Hàm kích hoạt sigmoid} & Sigmoid activation function \\
\textbf{Hàm mất mát đa dạng kênh} & Channel-diverse loss \\
\textbf{Hàm mất mát phân loại} & Classification loss \\
\textbf{Hàm mất mát phân tách} & Separation loss \\
\textbf{Hệ thống giám sát người lái xe} & Driver monitoring system \\
\textbf{Hệ thống phân tích biểu cảm khuôn mặt tự động} & Automated facial expression analysis \\
\textbf{Học bán giám sát} & Semi-supervised learning \\
\textbf{Học máy} & Machine learning \\
\textbf{Học phân phối nhãn} & Label distribution learning \\
\textbf{Học quá khớp} & Overfit / Overfitting \\
\textbf{Học vẹt} & Rote learning \\
\textbf{Huấn luyện từ đầu} & From scratch \\
\textbf{Kênh đặc trưng} & Feature channel \\
\textbf{Khả năng tổng quát hóa} & Generalizability / Generalization \\
\textbf{Khả năng tổng quát hóa chéo miền} & Cross-domain generalization \\
\textbf{Khả năng tổng quát hóa zero-shot} & Zero-shot generalization \\
\textbf{Khối chuyển đổi có khả năng học} & Learnable projection head \\
\textbf{Kiến trúc nhánh kép} & Dual-branch architecture \\
\textbf{Lô nhỏ} & Mini-batch \\
\textbf{Luồng quang học} & Optical flow \\
\textbf{Lớp kết nối đầy đủ} & Fully connected layer \\
\textbf{Mạng giáo viên} & Teacher network \\
\textbf{Mạng học sinh} & Student network \\
\textbf{Mạng nơ-ron} & Neural network \\
\textbf{Mạng nơ-ron đa tầng} & Multi-layer perceptron \\
\textbf{Mạng nơ-ron tích chập} & Convolutional neural network \\
\textbf{Mạng xương sống} & Backbone \\
\textbf{Mạng xương sống tiền huấn luyện} & Pre-trained backbone \\
\textbf{Mất cân bằng dữ liệu} & Data imbalance \\
\textbf{Mặt nạ} & Mask \\
\textbf{Mặt nạ sigmoid} & Sigmoid mask \\
\textbf{Miền dữ liệu} & Domain \\
\textbf{Miền dữ liệu đích (Miền đích)} & Target domain \\
\textbf{Miền dữ liệu nguồn (Miền nguồn)} & Source domain \\
\textbf{Mô-đun phân tách kênh} & Channel-separation module \\
\textbf{Mô hình diện mạo chủ động} & Active appearance model \\
\textbf{Mô hình hóa dựa trên tác tử} & Agent-based modeling \\
\textbf{Mô hình ngôn ngữ - hình ảnh} & Vision-language model \\
\textbf{Nghiên cứu cắt bỏ} & Ablation study \\
\textbf{Nhãn giả} & Pseudo-label \\
\textbf{Nhãn nhiễu} & Noisy labels \\
\textbf{Nhãn thực tế} & Ground truth \\
\textbf{Nhận dạng cảm xúc khuôn mặt} & Facial Expression Recognition \\
\textbf{Nhận dạng cảm xúc khuôn mặt có khả năng tổng quát hóa} & Generalizable Facial Expression Recognition \\
\textbf{Nhận dạng cảm xúc khuôn mặt động} & Dynamic facial expression recognition \\
\textbf{Nhận dạng cảm xúc khuôn mặt tĩnh} & Static facial expression recognition \\
\textbf{Nhận dạng khuôn mặt} & Face recognition \\
\textbf{Nội dung thiên lệch phương Tây} & Western-centric content \\
\textbf{Phản hồi thấu cảm} & Empathic expression \\
\textbf{Phân phối nhãn} & Label distribution \\
\textbf{Phân tách kênh} & Channel-separation \\
\textbf{Phép toán gộp cực đại} & Max pooling \\
\textbf{Phép toán trung bình} & Mean operation \\
\textbf{Quá trình hiệu chỉnh} & Calibration \\
\textbf{Quy luật} & Patterns \\
\textbf{Siêu tham số} & Hyperparameter \\
\textbf{Sự đánh đổi} & Trade-off \\
\textbf{Sự khác biệt cá nhân} & Individual differences \\
\textbf{Sự khác biệt về miền dữ liệu} & Domain gap \\
\textbf{Sự không khớp chiều không gian đặc trưng} & Feature dimension mismatch \\
\textbf{Tập dữ liệu kiểm tra chưa từng thấy} & Unseen test set \\
\textbf{Thế giới thực} & Real-world \\
\textbf{Thích ứng miền} & Domain adaptation \\
\textbf{Thiên lệch văn hóa trong nhận dạng cảm xúc} & Cross-cultural facial expression recognition bias \\
\textbf{Thu thập khuôn mặt} & Face acquisition \\
\textbf{Tinh chỉnh} & Fine-tune \\
\textbf{Tổng quát hóa liền mạch} & Generalize seamlessly \\
\textbf{Trích xuất và biểu diễn dữ liệu khuôn mặt} & Facial data extraction and representation \\
\textbf{Triển khai thực tế} & Real-world deployment \\
\textbf{Trợ lý ảo} & Virtual human \\
\textbf{Tương tác người - máy} & Human-robot interaction \\
\textbf{Xác suất phân loại (Hệ số xác suất)} & Logit \\

\end{longtable}

\pagebreak
\section*{Danh sách các ký hiệu và chữ viết tắt}
\addcontentsline{toc}{section}{Danh sách các ký hiệu và chữ viết tắt}

\renewcommand{\arraystretch}{1.3} 

\begin{longtable}{p{3.5cm} p{11.5cm}}
%─────────────────────────────────────────────────────────────────────────────
% PHẦN 1 — CHỮ VIẾT TẮT
%─────────────────────────────────────────────────────────────────────────────
\rowcolor{headerblue!25}
\multicolumn{2}{l}{\bfseries\color{headerblue}
    \quad Chữ viết tắt} \\
\endfirsthead

\textbf{AAM} & Mô hình diện mạo chủ động (Active Appearance Model) \\
\textbf{AFEA} & Hệ thống phân tích biểu cảm khuôn mặt tự động (Automated Facial Expression Analysis) \\
\textbf{API} & Giao diện lập trình ứng dụng (Application Programming Interface) \\
\textbf{AU} & Đơn vị hành động cơ mặt (Action Unit) \\
\textbf{AUC} & Diện tích dưới đường cong (Area Under the Curve) \\
\textbf{CAFE} & Phương pháp mô phỏng nhận thức con người cho nhận dạng cảm xúc khuôn mặt (Cognition of humAn for Facial Expression) \\
\textbf{CLIP} & Mô hình tiền huấn luyện ngôn ngữ và hình ảnh đối lập (Contrastive Language-Image Pre-training) \\
\textbf{DFER} & Nhận dạng cảm xúc khuôn mặt động (Dynamic Facial Expression Recognition) \\
\textbf{DMS} & Hệ thống giám sát người lái xe (Driver Monitoring System) \\
\textbf{EAC} & Phương pháp học nhất quán dựa trên xóa vùng ảnh (Erasure-based Consistency) \\
\textbf{FACS} & Hệ thống mã hóa hành động khuôn mặt (Facial Action Coding System) \\
\textbf{FC} & Lớp kết nối đầy đủ (Fully Connected) \\
\textbf{FER} & Nhận dạng cảm xúc khuôn mặt (Facial Expression Recognition) \\
\textbf{FR} & Nhận dạng khuôn mặt (Face Recognition) \\
\textbf{GFER} & Nhận dạng cảm xúc khuôn mặt có khả năng tổng quát hóa (Generalizable Facial Expression Recognition) \\
\textbf{GRU} & Khối nơ-ron có cổng (Gated Recurrent Unit) \\
\textbf{HMM} & Mô hình Markov ẩn (Hidden Markov Model) \\
\textbf{HRI} & Tương tác người - máy (Human-Robot Interaction) \\
\textbf{LSTM} & Mạng bộ nhớ dài - ngắn hạn (Long Short-Term Memory) \\
\textbf{MA} & Độ chính xác trung bình (Mean Accuracy) \\
\textbf{MLR} & Hồi quy logistic đa thức (Multinomial Logistic Regression) \\
\textbf{OA} & Độ chính xác tổng thể (Overall Accuracy) \\
\textbf{OFER} & Phương pháp nhận dạng biểu cảm khuôn mặt chịu lỗi (Occultation-aware FER) \\
\textbf{RUL} & Phương pháp học sự không chắc chắn (Relative Uncertainty Learning) \\
\textbf{SCN} & Mạng lọc nhiễu tự phục hồi (Self-Cure Network) \\
\textbf{SOTA} & Các phương pháp tiên tiến nhất (State-Of-The-Art) \\
\textbf{SSL} & Học bán giám sát (Semi-Supervised Learning) \\
\textbf{SVM} & Máy học vector hỗ trợ (Support Vector Machine) \\

\multicolumn{2}{c}{} \\

%─────────────────────────────────────────────────────────────────────────────
% PHẦN 2 — KÝ HIỆU TOÁN HỌC
%─────────────────────────────────────────────────────────────────────────────
\rowcolor{headerblue!25}
\multicolumn{2}{l}{\bfseries\color{headerblue}
    \quad Ký hiệu toán học} \\
    

$C$ & Số chiều của khối đặc trưng \\
$c$ & Hằng số chuẩn hóa \\
$D_{real}$ & Tập dữ liệu thực tế (miền dữ liệu đích) \\
$D_{train}$ & Tập dữ liệu huấn luyện (miền dữ liệu nguồn) \\
$F$ & Khối đặc trưng khuôn mặt tổng quát được trích xuất từ mô hình ngôn ngữ - hình ảnh lớn \\
$\tilde{F}$ & Khối đặc trưng biểu cảm cuối cùng sau khi được chọn lọc qua mặt nạ \\
$F^d$ & Ma trận các giá trị xác suất phân loại dự đoán \\
$\tilde{F}_{max}$ & Ma trận chứa các giá trị cực đại trích xuất dọc theo chiều kênh \\
$f$ & Ma trận đặc trưng biểu cảm thô trích xuất từ mạng học sâu khởi tạo mới \\
$L$ & Số lượng các lớp cảm xúc cơ bản đầu ra \\
$l_{cls}$ & Hàm mất mát phân loại \\
$l_{div}$ & Hàm mất mát đa dạng kênh \\
$l_{sep}$ & Hàm mất mát phân tách \\
$l_{train}$ & Tổng hàm mất mát huấn luyện \\
$M$ & Ma trận mặt nạ đặc trưng ban đầu \\
$M_{drop}$ & Ma trận rớt kênh ngẫu nhiên \\
$M_s$ & Ma trận mặt nạ sau khi đi qua hàm kích hoạt sigmoid \\
$N$ & Số lượng mẫu ảnh trong một cấu trúc dữ liệu tính toán (ví dụ: kích thước batch) \\
$x_i, \tilde{x}_i$ & Hình ảnh khuôn mặt đầu vào \\
$y_i$ & Nhãn cảm xúc thực tế tương ứng \\
$\lambda, \beta$ & Các siêu tham số cân bằng hàm mất mát \\

\end{longtable}

\renewcommand{\arraystretch}{1.0}